% EMNLP 2025 submission: Code2LoRA
\documentclass[11pt]{article}
\usepackage{EMNLP2025}
\usepackage{times}
\usepackage{latexsym}
\usepackage{amsmath}
\usepackage{amssymb}
\usepackage{graphicx}
\usepackage{booktabs}
\usepackage{multirow}
\usepackage{xcolor}
\usepackage{subcaption}
\usepackage{algorithm}
\usepackage{algorithmic}
\usepackage{url}
\usepackage{enumitem}
\usepackage{hyperref}
\usepackage{colortbl}
\usepackage{pifont}

\newcommand{\method}{Code2LoRA}
\newcommand{\benchmark}{RepoPeftBench}
\newcommand{\todo}[1]{\textcolor{red}{\textbf{[TODO: #1]}}}
\newcommand{\todoinline}[1]{\textcolor{red}{#1}}

\title{\method{}: Repository-Aware Adapter Generation\\via Composable Hypernetworks}

\author{
  Author 1 \and Author 2 \\
  Affiliation \\
  \texttt{\{author1, author2\}@institution.edu}
}

\begin{document}
\maketitle

% ============================================================
\begin{abstract}
% ============================================================

Large language models for code benefit from repository-level context, yet incorporating full codebases at inference time is expensive and bounded by context windows.
We propose \method{}, a hypernetwork that generates lightweight LoRA adapters conditioned on repository code embeddings.
Given a new repository, \method{} produces repo-specific adapter weights in a single forward pass---requiring zero additional inference tokens and no per-repo retraining.
We introduce two hypernetwork architectures:
(1) a \emph{direct-projection} variant that maps repository embeddings to full LoRA matrices via per-module-type output heads, and
(2) a \emph{shared-basis} variant (PAW-style) that generates per-layer mixing coefficients over learnable basis vectors, reducing hypernetwork parameters by $\sim$8.5$\times$.
We further propose composable LoRA generation, where file-level adapters are independently generated and combined, enabling incremental adaptation as codebases evolve.
We evaluate on \benchmark{}, a new benchmark of 560 Python repositories with 73K assertion completion tasks spanning cross-repo (CR) and in-repo (IR) evaluation settings.
\method{} achieves 63.8\% exact match on CR evaluation, substantially outperforming pretrained baselines (45.7\%), full fine-tuning (51.4\%), retrieval-augmented generation (39.9\%), and in-context learning (42.2\%).
\todo{Update numbers after all experiments complete.}

\end{abstract}

% ============================================================
\section{Introduction}
\label{sec:intro}
% ============================================================

Code completion and understanding tasks benefit significantly from repository-level context---knowledge about project structure, coding conventions, utility functions, and domain-specific patterns.
While retrieval-augmented generation (RAG) can provide some of this context by prepending relevant code chunks to the input, this approach has fundamental limitations: it increases inference cost linearly with context size, is bounded by the model's context window, and treats all queries identically regardless of the repository.
In-context learning (ICL) offers another route by providing demonstration examples, but is similarly constrained by context limits and adds latency at every inference step.

We propose an alternative paradigm: rather than retrieving raw code at inference time, we \emph{distill} repository knowledge into parameter-efficient adapters.
\method{} is a hypernetwork~\cite{ha2017hypernetworks} that generates LoRA~\cite{hu2022lora} adapter weights conditioned on repository code embeddings.
Given a code repository, the hypernetwork produces repo-specific LoRA weights in a single forward pass through a small neural network, specializing a frozen language model for that repository.

Our approach offers three key advantages over existing methods:
\begin{enumerate}[nosep]
    \item \textbf{Zero-token overhead}: Generated adapters modify model weights directly, adding zero inference tokens---unlike RAG, which requires processing retrieved chunks for every query.
    \item \textbf{Composability}: We introduce file-level LoRA generation, where adapters for individual source files are independently generated and composed. When a file changes, only its adapter needs regeneration.
    \item \textbf{Cross-repo generalization}: The hypernetwork generalizes to unseen repositories, producing effective adapters for codebases never seen during training---a capability that per-repo LoRA training inherently lacks.
\end{enumerate}

We investigate two hypernetwork architectures:
a \emph{direct-projection} design that maps repository embeddings to full LoRA matrices through per-module-type output heads with learnable scaling,
and a \emph{shared-basis} design (PAW-style) that maintains learnable basis vectors per module type and generates only per-layer mixing coefficients, reducing the hypernetwork parameter count by $\sim$8.5$\times$ while maintaining comparable performance.

We evaluate \method{} on \benchmark{}, a new benchmark we construct from 560 Python repositories containing 73,267 assertion completion tasks.
We define two evaluation settings: \emph{cross-repo} (CR), where the model encounters entirely new repositories at test time, and \emph{in-repo} (IR), where the model has seen other test functions from the same repository during training.
Against a comprehensive set of seven baselines---pretrained, RAG, ICL, oracle context, full fine-tuning, single LoRA, and per-repo LoRA---\method{} achieves the best performance on cross-repo evaluation while requiring no additional inference tokens.
\todo{Verify final claim after all experiments.}

Our contributions are:
\begin{enumerate}[nosep]
    \item We propose \method{}, the first hypernetwork for generating repository-conditioned LoRA adapters for code LLMs, enabling instant specialization for unseen repositories.
    \item We introduce two hypernetwork architectures---direct-projection and shared-basis---and analyze their trade-offs in parameter efficiency and generation quality.
    \item We construct \benchmark{}, a benchmark of 560 Python repositories with carefully curated assertion completion tasks, structured prefixes, and dual evaluation protocols.
    \item We conduct comprehensive experiments against seven baselines with and without oracle context, demonstrating that \method{} substantially outperforms all context-based methods while maintaining efficiency comparable to no-context baselines.
\end{enumerate}

% ============================================================
\section{Related Work}
\label{sec:related}
% ============================================================

\paragraph{Parameter-Efficient Fine-Tuning.}
LoRA~\cite{hu2022lora} enables efficient model adaptation through low-rank decomposition of weight updates.
Extensions include LoRA composition~\cite{zhang2023composing}, weight merging~\cite{yadav2024ties}, and multi-LoRA routing~\cite{huang2024lorahub}.
\todo{Add QLoRA citation, DoRA citation.}
These approaches treat adapters as static artifacts trained for individual tasks.
Our work instead dynamically \emph{generates} adapters conditioned on input data, enabling adaptation to unseen domains without any training.

\paragraph{Hypernetworks.}
Hypernetworks~\cite{ha2017hypernetworks} generate the parameters of a target network from a conditioning signal.
Recent applications to language models include generating LoRA weights conditioned on task descriptions~\cite{phang2023hypertuning} and document content~\cite{von2024language}.
\todo{Add HyperTuning citation details, check for concurrent work in code domain.}
Our work applies hypernetworks specifically to code repositories, introducing composition mechanisms for file-level adapter generation and evaluating with a purpose-built code benchmark.

\paragraph{Repository-Level Code Understanding.}
Prior work on incorporating repository context includes RepoFusion~\cite{shrivastava2023repofusion}, which trains with cross-file context, RepoCoder~\cite{zhang2023repocoder}, which iteratively retrieves and generates, and CrossCodeEval~\cite{ding2024crosscodeeval}, a cross-file evaluation benchmark.
\todo{Add RepoFormer, CoCoMIC, R2C2-Coder citations.}
These approaches typically incorporate repository knowledge through the input (retrieval or cross-file context).
In contrast, \method{} distills repository knowledge into model \emph{parameters}, avoiding context-window limitations and runtime retrieval costs.

\paragraph{Code LLMs.}
Modern code language models including CodeLlama~\cite{roziere2023code}, StarCoder~\cite{li2023starcoder}, and Qwen2.5-Coder~\cite{hui2024qwen25coder} achieve strong performance across code tasks.
We use Qwen2.5-Coder-1.5B as our base model due to its strong performance at small scale (1.5B parameters), extensive pretraining on code data, and support for fill-in-the-middle generation.
\todo{Add DeepSeek-Coder citation, verify no newer base model should be used.}

% ============================================================
\section{Method}
\label{sec:method}
% ============================================================

\subsection{Overview}

\method{} consists of three components (Figure~\ref{fig:architecture}):
(1) a \textbf{repository encoder} that produces dense embeddings from source code files,
(2) a \textbf{hypernetwork} that generates LoRA adapter weights from these embeddings, and
(3) a \textbf{frozen code LLM} that receives the generated adapters and performs inference.
At training time, only the hypernetwork parameters are optimized through the language modeling loss; the base LLM remains entirely frozen.

\begin{figure}[t]
\centering
\todo{Architecture diagram: repo files $\to$ encoder $\to$ embedding $\to$ hypernetwork $\to$ LoRA weights $\to$ frozen LLM + LoRA $\to$ output}
\caption{Overview of the \method{} architecture. Source files are encoded into a repository embedding, which conditions a hypernetwork that generates LoRA adapter weights. The frozen base LLM is specialized via these generated adapters without any retraining.}
\label{fig:architecture}
\end{figure}

\subsection{Repository Encoding}
\label{sec:encoding}

We encode repository source files using Qwen3-Embedding-0.6B, a 0.6B-parameter dense embedding model pretrained on diverse text including code.
Non-test source files are collected from each repository, excluding common non-informative directories (e.g., \texttt{.git}, \texttt{\_\_pycache\_\_}, \texttt{venv}).

Each source file is chunked into windows of up to 4,096 tokens with 512-token overlap.
File-level embeddings are obtained via mean pooling over chunk embeddings.
Files are assigned importance weights based on three heuristics:
(i) \emph{content distinctiveness}, measured as cosine distance from the repository mean embedding;
(ii) \emph{file size}, with larger files receiving higher weight;
(iii) \emph{path heuristics}, upweighting files in \texttt{src/}, \texttt{core/}, \texttt{lib/} directories and downweighting boilerplate files such as \texttt{\_\_init\_\_.py}, \texttt{setup.py}, and \texttt{version.py}.

The repository embedding is formed by concatenating a weighted mean and max-pooling over file embeddings:
\begin{equation}
    \mathbf{e}_\text{repo} = \left[\sum_i w_i \mathbf{f}_i \;;\; \max_i \mathbf{f}_i\right] \in \mathbb{R}^{2d}
\label{eq:repo_emb}
\end{equation}
where $\mathbf{f}_i \in \mathbb{R}^{d}$ is the embedding of file $i$, $w_i$ are normalized importance weights, and $d = 1024$ is the embedding dimension of Qwen3-Embedding.
The concatenation of mean and max pooling captures both the average ``flavor'' of a repository and its most distinctive features.
\todo{Report embedding statistics: mean files/repo, distribution of embedding norms.}

\subsection{Hypernetwork Architectures}
\label{sec:hypernetwork}

We investigate two hypernetwork architectures that trade off between parameter efficiency and representational capacity.

\subsubsection{Direct-Projection Hypernetwork}
\label{sec:direct_proj}

The direct-projection hypernetwork $H_\theta$ maps a repository embedding to LoRA parameters through a shared trunk and per-module-type output heads.
The shared trunk is a 2-layer MLP:
\begin{equation}
    \mathbf{h} = \text{MLP}(\mathbf{e}_\text{repo}) = \sigma(\mathbf{W}_2 \cdot \sigma(\mathbf{W}_1 \cdot \mathbf{e}_\text{repo}))
\end{equation}
where $\sigma$ denotes the GELU activation.
The trunk output is L2-normalized and rescaled by $\sqrt{d_h}$ to stabilize gradient flow.

For each module type $t$ (e.g., \texttt{q\_proj}, \texttt{k\_proj}, \texttt{v\_proj}, \texttt{o\_proj}, \texttt{gate\_proj}, \texttt{up\_proj}, \texttt{down\_proj} in a transformer), dedicated output heads generate the LoRA $A$ and $B$ matrices:
\begin{align}
    \mathbf{A}_t &= \tanh(\text{Head}^A_t(\mathbf{h})) \cdot \exp(s^A_t) \\
    \mathbf{B}_t &= \tanh(\text{Head}^B_t(\mathbf{h})) \cdot \exp(s^B_t)
\end{align}
where $s^A_t, s^B_t$ are learnable log-scale parameters initialized to $-3.5$ (corresponding to an initial scale of $\approx 0.03$), ensuring small initial adapter magnitudes.
The $\tanh$ activation bounds raw values to $[-1, 1]$, and the exponential scale (clamped to $[10^{-5}, 0.3]$) provides a learnable magnitude control.

The same LoRA matrices are shared across all layers of the same module type.
For LoRA rank $r = 16$ and the 7 module types in Qwen2.5-Coder's 28-layer architecture, this produces 196 $(A, B)$ matrix pairs.
These are injected by directly modifying the linear layer weights:
\begin{equation}
    \mathbf{W}' = \mathbf{W} + \mathbf{B}_t \mathbf{A}_t
\end{equation}

\paragraph{Parameter count.}
With input dimension $2d = 2048$, hidden dimension 512, rank 16, and 7 module types with varying dimensions, the direct-projection hypernetwork has approximately \todo{compute exact number, $\sim$50M} parameters.

\subsubsection{Shared-Basis Hypernetwork (PAW-Style)}
\label{sec:paw}

The shared-basis variant addresses the parameter inefficiency of direct projection by factoring LoRA generation through learnable basis vectors.

\paragraph{Architecture.}
For each module type $m$ (e.g., \texttt{q\_proj}), the hypernetwork maintains a set of $N$ learnable basis matrices:
\begin{align}
    \mathbf{A}^m_\text{bases} &\in \mathbb{R}^{N \times r \times d_\text{in}} \quad \text{(normal init, } \sigma = 0.02\text{)} \\
    \mathbf{B}^m_\text{bases} &\in \mathbb{R}^{N \times d_\text{out} \times r} \quad \text{(zero init, per LoRA convention)}
\end{align}

A shared trunk with residual MLP blocks processes the repository embedding:
\begin{equation}
    \mathbf{h} = \text{ResidualTrunk}(\mathbf{e}_\text{repo})
\end{equation}
where each residual block applies LayerNorm $\to$ Linear $\to$ GELU $\to$ Linear $\to$ residual add.

A coefficient head maps the trunk output to per-layer, per-module mixing weights:
\begin{equation}
    \boldsymbol{\alpha} = \text{CoeffHead}(\mathbf{h}) \in \mathbb{R}^{L \times M \times N \times 2}
\end{equation}
where $L$ is the number of transformer layers, $M$ the number of module types, and $N$ the number of bases.
The final per-layer LoRA matrices are computed via linear combination:
\begin{align}
    \mathbf{A}^{(l,m)} &= \sum_{n=1}^{N} \alpha^A_{l,m,n} \cdot \mathbf{A}^m_{\text{bases},n} \\
    \mathbf{B}^{(l,m)} &= \sum_{n=1}^{N} \alpha^B_{l,m,n} \cdot \mathbf{B}^m_{\text{bases},n}
\end{align}

\paragraph{Hook-based injection.}
Unlike the direct-projection variant that modifies weights in-place, the shared-basis variant applies LoRA through forward hooks:
\begin{equation}
    \mathbf{y}' = \mathbf{y} + (\mathbf{x} \mathbf{A}^T \mathbf{B}^T) \cdot \frac{\alpha}{r}
\end{equation}
where $\mathbf{x}$ is the input, $\mathbf{y}$ is the original output, and $\alpha/r$ is the LoRA scaling factor.
Hooks are registered before the forward pass and removed after, enabling clean batch-level adapter switching.

\paragraph{Parameter count.}
The shared-basis design requires only $N \times (r \times d_\text{in} + d_\text{out} \times r)$ parameters for basis storage per module type, plus $L \times M \times N \times 2$ mixing coefficients per input.
With $N = 16$ bases, this achieves $\sim$8.5$\times$ fewer parameters than the direct-projection variant.
\todo{Report exact parameter counts for both variants.}

\paragraph{Comparison.}
Table~\ref{tab:arch_comparison} summarizes the architectural differences.
The direct-projection variant generates layer-\emph{shared} LoRA weights (same matrices applied to all layers of a given type), while the shared-basis variant generates \emph{per-layer} weights through mixing coefficients, providing finer-grained control.

\begin{table}[t]
\centering
\small
\begin{tabular}{lcc}
\toprule
Property & Direct-Proj. & Shared-Basis \\
\midrule
Per-layer weights & \ding{55} & \ding{51} \\
Basis sharing & \ding{55} & \ding{51} \\
Trunk depth & 2 (MLP) & 2+ (Residual) \\
Injection & Weight mod. & Fwd hooks \\
Params ($\sim$) & \todo{50M} & \todo{6M} \\
\bottomrule
\end{tabular}
\caption{Architectural comparison of the two hypernetwork variants.}
\label{tab:arch_comparison}
\end{table}

\subsection{Composable LoRA Generation}
\label{sec:composable}

Beyond monolithic repository embeddings, we propose file-level LoRA generation and composition.
For each source file $f_i$ in a repository, the hypernetwork generates a separate adapter:
\begin{equation}
    (\mathbf{A}_i, \mathbf{B}_i) = H_\theta(\mathbf{f}_i)
\end{equation}

These file-level adapters are composed into a repository-level adapter using one of three strategies:

\paragraph{Additive.} Simple averaging: $\mathbf{A} = \frac{1}{K}\sum_{i=1}^K \mathbf{A}_i$.
This provides a uniform combination of all file-level knowledge.

\paragraph{Weighted.} Learned importance: $\mathbf{A} = \sum_{i=1}^K \text{softmax}(g(\mathbf{f}_i))_i \cdot \mathbf{A}_i$ where $g$ is a small weight predictor trained end-to-end.
This allows the hypernetwork to attend more strongly to informative files.

\paragraph{Gated.} Query-conditioned: Composition weights depend on both file embeddings and the input query, enabling dynamic selection of relevant file knowledge per test case.

The composable design enables \emph{incremental adaptation}: when a single file in a repository changes, only its LoRA needs regeneration, and the composed adapter is updated in $O(1)$ time.

\todo{Report composition results. If not ready, describe as ``we leave composable evaluation to future work'' and focus on repo-level.}

\subsection{Training}
\label{sec:training}

The hypernetwork is trained end-to-end via the standard language modeling loss on assertion completion tasks.
The base LLM remains entirely frozen; only hypernetwork parameters $\theta$ are optimized:
\begin{equation}
    \mathcal{L}(\theta) = -\sum_{(x,y,r) \in \mathcal{D}} \log p_{\text{LLM} + \text{LoRA}(H_\theta(\mathbf{e}_r))}(y \mid x)
\label{eq:training_loss}
\end{equation}
where $x$ is the test prefix (structured context before the assertion's expected value), $y$ is the assertion target to be completed, $\mathbf{e}_r$ is the embedding of repository $r$, and $\text{LoRA}(H_\theta(\mathbf{e}_r))$ denotes the adapter weights generated by the hypernetwork.

\paragraph{Sampling strategy.}
Training batches are constructed by sampling a repository, then sampling a QnA pair from that repository.
This ensures the hypernetwork sees diverse repositories across training and does not overfit to repositories with many examples.

\paragraph{Oracle context augmentation.}
Optionally, the test prefix $x$ can be augmented with import-resolved oracle context (Section~\ref{sec:oracle}).
When oracle context is available, it is prepended to the prefix, providing the hypernetwork with both the repository embedding (through LoRA) and explicit function definitions (through the input).
We train and evaluate with and without oracle augmentation to disentangle the effects of parametric adaptation and context injection.

% ============================================================
\section{\benchmark{}: A Repository-Level Code Benchmark}
\label{sec:benchmark}
% ============================================================

We construct \benchmark{}, a benchmark for evaluating repository-level code understanding through assertion completion.

\subsection{Repository Collection}

We collect 560 Python repositories from GitHub, filtered by:
(i) use of \texttt{pytest} or \texttt{unittest} testing frameworks;
(ii) at least 6 GitHub stars;
(iii) recent activity (updated within the past year);
(iv) permissive open-source licenses.
For each repository, we separate source files into \emph{test files} (used for QnA extraction) and \emph{non-test files} (used for repository encoding).
Test files are identified by path heuristics (\texttt{test\_*.py}, \texttt{*\_test.py}, files in \texttt{tests/} directories) and moved to a dedicated directory.

\subsection{QnA Pair Extraction}
\label{sec:qna_extraction}

We extract assertion completion tasks using AST (Abstract Syntax Tree) parsing.
For each test file, we:

\begin{enumerate}[nosep]
    \item Parse the file into an AST and identify assertion nodes of five types: bare \texttt{assert} statements, \texttt{self.assert*} method calls, \texttt{pytest.raises} context managers, \texttt{pytest.approx} expressions, and NumPy-style \texttt{assert\_*} calls.
    \item For each assertion, extract the \emph{target}---the expression being verified---as the completion target.
    \item Construct a \emph{structured prefix} consisting of: (a) import statements from the test file, (b) the enclosing class definition (if any), (c) helper methods (setUp/tearDown), and (d) the test function body up to the assertion target.
    \item Apply quality filters: remove targets starting with commas (malformed AST segmentation), targets outside function bodies, empty or whitespace-only targets, and duplicate targets within the same function.
\end{enumerate}

\paragraph{Difficulty classification.}
Each QnA pair is assigned a difficulty level based on the target complexity:
\emph{simple} (single value or short expression, $\leq 30$ characters),
\emph{medium} (function call, attribute access, or moderate expression, 30--80 characters), and
\emph{hard} (nested calls, complex expressions, or multi-part assertions, $>$80 characters).
We ensure balanced difficulty distribution within each split.

\subsection{Evaluation Protocol}

We define two complementary evaluation settings:

\paragraph{Cross-Repo (CR).}
58 repositories are entirely held out during training.
This evaluates the hypernetwork's ability to generalize to completely unseen codebases---the primary use case for \method{}.
These repositories have no overlap with training data at the repository level.

\paragraph{In-Repo (IR).}
Test assertions come from 447 repositories seen during training, but from different test \emph{functions} than those used for training.
This evaluates within-repository generalization---whether adaptation to a repository transfers to unseen test cases from the same codebase.
Per-repo LoRA is applicable only in this setting.

\subsection{Oracle Context}
\label{sec:oracle}

For a subset of QnA pairs, we construct import-resolved oracle context.
For each test file's prefix, we:
(i) parse all import statements using AST analysis with fallback regex parsing for syntax errors;
(ii) resolve imported modules to source files within the repository, trying multiple source roots (\texttt{src/}, \texttt{lib/}, package directories);
(iii) identify names actually used in the test prefix;
(iv) extract the specific function and class definitions being tested.

The extracted definitions are concatenated (up to a token budget) and prepended to the test prefix.
This provides an upper bound on the information available through import analysis.
\todo{Report oracle coverage: $\sim$87.5\% of pairs have non-empty oracle context.}

\subsection{Dataset Statistics}

\begin{table}[t]
\centering
\small
\begin{tabular}{lrrr}
\toprule
Split & Repos & QnA Pairs & Pairs/Repo \\
\midrule
Train & 447 & 46,313 & 103.6 \\
IR Val & 447 & 5,679 & 12.7 \\
IR Test & 447 & 5,222 & 11.7 \\
CR Val & 55 & 7,609 & 138.3 \\
CR Test & 58 & 6,414 & 110.6 \\
\midrule
Total & 560 & 71,237 & 127.2 \\
\bottomrule
\end{tabular}
\caption{\benchmark{} dataset statistics. Counts reflect pairs after quality filtering (comma-leading target removal). \todo{Verify final counts; total was 73,267 pre-filter.}}
\label{tab:dataset_stats}
\end{table}

\begin{table}[t]
\centering
\small
\begin{tabular}{lrrr}
\toprule
Property & Train & CR Test & IR Test \\
\midrule
Assertion types & & & \\
\quad \texttt{assert} & \todo{X\%} & \todo{X\%} & \todo{X\%} \\
\quad \texttt{self.assert*} & \todo{X\%} & \todo{X\%} & \todo{X\%} \\
\quad \texttt{pytest.*} & \todo{X\%} & \todo{X\%} & \todo{X\%} \\
\midrule
Difficulty & & & \\
\quad Simple & \todo{X\%} & \todo{X\%} & \todo{X\%} \\
\quad Medium & \todo{X\%} & \todo{X\%} & \todo{X\%} \\
\quad Hard & \todo{X\%} & \todo{X\%} & \todo{X\%} \\
\midrule
Avg.\ target length (chars) & \todo{X} & \todo{X} & \todo{X} \\
Avg.\ prefix length (tokens) & \todo{X} & \todo{X} & \todo{X} \\
Oracle coverage (\%) & \todo{X} & \todo{X} & \todo{X} \\
\bottomrule
\end{tabular}
\caption{Detailed dataset statistics across splits. \todo{Fill from validation script output.}}
\label{tab:dataset_details}
\end{table}

% ============================================================
\section{Experimental Setup}
\label{sec:setup}
% ============================================================

\subsection{Base Model and Embeddings}

\paragraph{Base LLM.}
Qwen2.5-Coder-1.5B~\cite{hui2024qwen25coder}, a 1.5B-parameter code language model.
The model is loaded in \texttt{bfloat16} precision and kept frozen during all hypernetwork training.
All baselines that involve fine-tuning also use this model.

\paragraph{Embedding model.}
Qwen3-Embedding-0.6B for repository encoding.
Each file is chunked into 4,096-token windows with 512-token overlap.
The same embedding model is used for RAG chunk indexing and ICL example retrieval.
\todo{Verify embedding model name and dimensions.}

\subsection{Hypernetwork Configuration}

\paragraph{Direct-projection (\method{}).}
Input dimension: 2,048 (concatenated mean+max of 1,024-dim embeddings).
Hidden dimension: 512.
LoRA rank: 16, alpha: 32.
Target modules: all 7 attention and MLP projection types across all 28 transformer layers (196 total injection points).
Trunk: 2-layer MLP with GELU activation and L2 normalization.

\paragraph{Shared-basis (\method{}-PAW).}
Same input and LoRA configuration.
Hidden dimension: 512.
Number of bases: $N = 16$.
Trunk depth: 2 (input projection + 1 residual MLP block).
Coefficient head maps to $L \times M \times N \times 2 = 28 \times 7 \times 16 \times 2 = 6{,}272$ mixing coefficients.

\paragraph{Training.}
AdamW optimizer with learning rate $1 \times 10^{-4}$, cosine schedule with warmup, 3 epochs, effective batch size 1 (gradient accumulation as needed).
Maximum sequence length: 8,192 tokens.
Training uses the SFT (Supervised Fine-Tuning) framework from TRL~\cite{vonwerra2022trl}.
Best checkpoint selected by validation loss on CR val.

\subsection{Baselines}
\label{sec:baselines}

We compare against seven baselines spanning inference-only and trained methods:

\paragraph{Pretrained.}
Qwen2.5-Coder-1.5B with no adaptation.
Generates completions directly from the test prefix with no additional context.

\paragraph{RAG ($k$=3).}
Non-test source files are pre-chunked into 512-token segments and embedded using Qwen3-Embedding.
At inference, the test prefix is embedded and top-$k$ chunks are retrieved by cosine similarity.
Retrieved chunks are prepended to the prefix.
We report $k=3$ (best performing); results for $k=5$ and $k=10$ are in the appendix.
\todo{Move k=5,10 results to appendix.}

\paragraph{ICL (3-shot).}
Few-shot examples from similar contexts.
For IR evaluation, examples are randomly sampled from same-repo training pairs.
For CR evaluation, examples are retrieved by embedding similarity between the test prefix and training prefixes across all repositories, ensuring diverse and relevant demonstrations.
\todo{Verify retrieval-based selection is used for final results.}

\paragraph{Oracle Context.}
Import-resolved source code definitions prepended to the prefix (Section~\ref{sec:oracle}).
This provides an upper bound on the information extractable from import analysis without any training.

\paragraph{Full Fine-Tuning (FFT).}
All parameters of Qwen2.5-Coder-1.5B are fine-tuned on the training set.
Training uses \texttt{bfloat16} precision, AdamW optimizer, learning rate $2 \times 10^{-5}$, 3 epochs, maximum sequence length 2,048 tokens.
Best checkpoint selected by validation loss.

\paragraph{Single LoRA (sLoRA).}
One LoRA adapter (rank 16, alpha 32) trained on all training data.
Training uses \texttt{bfloat16} precision, AdamW optimizer, 3 epochs, maximum sequence length 2,048 tokens.

\paragraph{Per-Repo LoRA (pLoRA).}
Individual LoRA adapters trained separately for each repository using only that repository's training data.
Applicable only for IR evaluation (adapters cannot be created for unseen CR repositories).
We train on 10 representative repositories.
\todo{Expand to all IR repos or report subset. Current: 10 repos, n=99 test pairs.}

\subsection{Evaluation Metrics}
\label{sec:metrics}

We report three complementary metrics:

\paragraph{Exact Match (EM).}
Binary metric: 1 if the predicted assertion target matches the reference after normalization (whitespace collapsing, trailing comma/colon removal).
We apply relaxed matching that handles common model overgeneration (e.g., predicting extra arguments after a comma when the reference has none).

\paragraph{Edit Similarity.}
The SequenceMatcher ratio~\cite{python_difflib} between prediction and reference, ranging from 0 (completely different) to 1 (identical).
Captures partial credit for near-correct predictions.

\paragraph{CodeBLEU.}
Code-aware BLEU~\cite{ren2020codebleu} that considers AST and data-flow structure in addition to n-gram overlap.
We use the \texttt{codebleu} package where available, with a fallback to tokenized BLEU computed on Python code tokens.
For exact matches, CodeBLEU is defined as 1.0 to avoid smoothing artifacts on short strings.
\todo{Add CodeBLEU citation.}

\paragraph{Inference configuration.}
All methods use identical inference settings: \texttt{max\_input\_tokens} = 16,384, \texttt{max\_new\_tokens} = 128, greedy decoding (temperature 0).
Post-processing strips FIM tokens, takes the first line (for single-line targets), and removes trailing comments.

% ============================================================
\section{Results}
\label{sec:results}
% ============================================================

\subsection{Main Results}

Table~\ref{tab:main_results} presents results across all methods on both evaluation settings.

\begin{table*}[t]
\centering
\small
\begin{tabular}{l@{\hskip 8pt}ccc@{\hskip 12pt}ccc}
\toprule
& \multicolumn{3}{c}{\textbf{Cross-Repo (CR Test)}} & \multicolumn{3}{c}{\textbf{In-Repo (IR Test)}} \\
\cmidrule(lr){2-4} \cmidrule(lr){5-7}
Method & EM (\%) & EditSim & CodeBLEU & EM (\%) & EditSim & CodeBLEU \\
\midrule
\multicolumn{7}{l}{\emph{Inference-only (no training)}} \\
\midrule
Pretrained & 45.7 & 0.597 & \todoinline{0.XX}\textsuperscript{\dag} & 46.8 & 0.615 & \todoinline{0.XX}\textsuperscript{\dag} \\
RAG ($k$=3) & 39.9 & 0.534 & 0.439 & 42.3 & 0.554 & 0.447 \\
ICL (3-shot) & 42.2 & 0.560 & 0.470 & 45.3 & 0.593 & 0.482 \\
Oracle Context & 47.5 & 0.609 & 0.484 & 48.7 & 0.618 & 0.486 \\
\midrule
\multicolumn{7}{l}{\emph{Trained (prefix $\to$ target)}} \\
\midrule
FFT & 51.4 & 0.367 & 0.355 & 55.9 & 0.389 & 0.379 \\
Single LoRA & 41.2 & 0.179 & 0.208 & 43.8 & 0.187 & 0.213 \\
Per-repo LoRA & --- & --- & --- & 57.6 & \todoinline{X.XX} & \todoinline{X.XX} \\
\textbf{\method{}} & \textbf{63.8} & \textbf{0.784} & \textbf{0.777} & \textbf{66.2} & \textbf{0.806} & \textbf{0.795} \\
\method{}-PAW & \todoinline{XX.X} & \todoinline{X.XX} & \todoinline{X.XX} & \todoinline{XX.X} & \todoinline{X.XX} & \todoinline{X.XX} \\
\midrule
\multicolumn{7}{l}{\emph{Trained (oracle + prefix $\to$ target)}} \\
\midrule
FFT + Oracle & \todoinline{XX.X} & \todoinline{X.XX} & \todoinline{X.XX} & \todoinline{XX.X} & \todoinline{X.XX} & \todoinline{X.XX} \\
sLoRA + Oracle & \todoinline{XX.X} & \todoinline{X.XX} & \todoinline{X.XX} & \todoinline{XX.X} & \todoinline{X.XX} & \todoinline{X.XX} \\
pLoRA + Oracle & --- & --- & --- & \todoinline{XX.X} & \todoinline{X.XX} & \todoinline{X.XX} \\
\textbf{\method{} + Oracle} & \todoinline{XX.X} & \todoinline{X.XX} & \todoinline{X.XX} & \todoinline{XX.X} & \todoinline{X.XX} & \todoinline{X.XX} \\
\method{}-PAW + Oracle & \todoinline{XX.X} & \todoinline{X.XX} & \todoinline{X.XX} & \todoinline{XX.X} & \todoinline{X.XX} & \todoinline{X.XX} \\
\bottomrule
\end{tabular}
\caption{Main results on \benchmark{}. EM = Exact Match (\%), EditSim = Edit Similarity, CodeBLEU. Best non-oracle results per column in \textbf{bold}. Per-repo LoRA is N/A for CR (no adapter for unseen repos). \textsuperscript{\dag}Pretrained CodeBLEU to be re-evaluated with \texttt{codebleu} package. \todo{Fill all missing cells. Re-run pretrained with CodeBLEU. Verify sLoRA numbers---current results use 4-bit quantized model; bf16 retrain in progress. pLoRA only on 10 repos (n=99); expand or caveat.}}
\label{tab:main_results}
\end{table*}

\paragraph{Key findings.}
\begin{enumerate}[nosep]
    \item \textbf{\method{} substantially outperforms all baselines on CR.}
    On cross-repo evaluation, \method{} achieves 63.8\% EM, outperforming the next best method (FFT, 51.4\%) by 12.4 absolute points and the pretrained baseline by 18.1 points.
    This demonstrates that repository embeddings encode substantial task-relevant knowledge.

    \item \textbf{Context-based methods underperform on assertion completion.}
    RAG (39.9\%) and ICL (42.2\%) perform \emph{below} the pretrained baseline (45.7\%) on CR.
    We hypothesize that retrieved chunks and few-shot examples consume context budget without providing directly useful information for assertion prediction, and may introduce distracting context.
    Oracle context (47.5\%) provides only a modest 1.8-point improvement, suggesting that even perfect import resolution offers limited benefit in the input space.

    \item \textbf{Parametric adaptation outperforms context injection.}
    All trained methods (except sLoRA) outperform context-based methods, confirming that adapting model parameters is more effective than input augmentation for repository-level tasks.
    \todo{Verify this claim holds after sLoRA bf16 retrain.}

    \item \textbf{IR consistently outperforms CR.}
    All methods show 2--4\% higher EM on IR vs.\ CR, reflecting the benefit of in-distribution test data from familiar repositories.

    \item \textbf{Per-repo LoRA is competitive on IR.}
    At 57.6\% EM on IR, per-repo LoRA demonstrates the value of repository-specific adaptation, but \method{} still outperforms it (66.2\%) while requiring no per-repo training.
    \todo{Verify pLoRA result is representative (only 10 repos, n=99). Consider caveating.}
\end{enumerate}

\paragraph{Discussion: RAG and ICL below pretrained.}
This counterintuitive result---context-based methods underperforming the pretrained baseline---warrants discussion.
Assertion completion targets are typically short expressions (variable names, function calls, constants) that depend heavily on \emph{local} context (the test function body) rather than retrieved repository content.
Prepending retrieved chunks or few-shot examples consumes context budget and can shift the model's attention away from the immediate test context.
We use a generous context window (16,384 tokens) to minimize truncation effects, but the fundamental issue is that retrieval noise outweighs the benefit of additional context for this particular task.
This motivates the \method{} approach: repository knowledge should be \emph{internalized} into model parameters rather than injected as context.
\todo{Run ablation varying context window size to support this claim.}

\subsection{Oracle Context Ablation}
\label{sec:oracle_ablation}

We evaluate whether oracle context at training time provides benefits that transfer to inference without oracle context (Table~\ref{tab:oracle_ablation}).

\begin{table}[t]
\centering
\small
\begin{tabular}{lcc}
\toprule
\multirow{2}{*}{Train Condition} & \multicolumn{2}{c}{Test EM (\%)} \\
\cmidrule(lr){2-3}
& w/o Oracle & w/ Oracle \\
\midrule
\multicolumn{3}{c}{\emph{CR Test}} \\
\midrule
\method{} & 63.8 & \todoinline{XX.X} \\
\method{} + Oracle & \todoinline{XX.X} & \todoinline{XX.X} \\
\midrule
\multicolumn{3}{c}{\emph{IR Test}} \\
\midrule
\method{} & 66.2 & \todoinline{XX.X} \\
\method{} + Oracle & \todoinline{XX.X} & \todoinline{XX.X} \\
\bottomrule
\end{tabular}
\caption{Oracle cross-evaluation ablation for \method{}. \todo{Fill after oracle-trained model completes. Key question: does ``train+oracle, test-oracle'' $\geq$ ``train-oracle, test-oracle''?}}
\label{tab:oracle_ablation}
\end{table}

\todo{Write analysis paragraph interpreting the oracle ablation results:
\begin{itemize}
    \item If train+oracle, test-oracle $\approx$ train-oracle, test-oracle $\to$ oracle training doesn't help without oracle at test (expected).
    \item If train+oracle, test-oracle $>$ train-oracle, test-oracle $\to$ model learned transferable patterns from oracle exposure (strong result).
    \item Discuss practical implications: when should practitioners use oracle context?
\end{itemize}}


\subsection{Composability Results}
\label{sec:composability}

\todo{If composable experiments are complete, report:
\begin{enumerate}
    \item Comparison of additive, weighted, and gated composition (Table or figure).
    \item Incremental adaptation: plot EM vs.\ number of files composed (monotonic improvement expected).
    \item Efficiency: time to update adapter when a single file changes.
\end{enumerate}
If not complete, replace this subsection with a sentence: ``We leave composable evaluation to future work.'' and remove from contributions list.}

\begin{figure}[t]
\centering
\todo{Incremental adaptation figure: x-axis = number of files, y-axis = EM. Show monotonic improvement.}
\caption{Incremental adaptation: performance improves as more file-level LoRAs are composed into the repository adapter.}
\label{fig:incremental}
\end{figure}

% ============================================================
\section{Analysis}
\label{sec:analysis}
% ============================================================

\subsection{Per-Repository Performance Distribution}

Figure~\ref{fig:per_repo_boxplot} shows the distribution of per-repo EM scores across IR-test repositories for different methods.

\begin{figure}[t]
\centering
\todo{Box/violin plot: per-repo EM for Pretrained, FFT, pLoRA, Code2LoRA on IR-test.
Expected patterns:
- pLoRA: high variance (small repos $\to$ few examples $\to$ unstable)
- Code2LoRA: lower variance (hypernetwork regularizes across repos)
- FFT/Pretrained: uniform (no per-repo signal)
Generate with: analysis/per\_repo\_boxplot.py}
\caption{Per-repository EM distribution on IR test. \method{} achieves consistently high performance with lower variance than per-repo LoRA, demonstrating the regularizing effect of the hypernetwork.}
\label{fig:per_repo_boxplot}
\end{figure}

\todo{Write analysis paragraph:
\begin{itemize}
    \item pLoRA high variance: small repos have few training examples, leading to unreliable adapters.
    \item Code2LoRA transfers knowledge across repos, stabilizing performance.
    \item Identify repos where Code2LoRA helps most (repos with unique coding patterns) vs.\ least (repos with generic code).
\end{itemize}}

\subsection{Scaling with Number of Training Repositories}

We train \method{} on subsets of 10, 20, 50, 100, 200, and all 447 training repositories and evaluate on CR test to measure the benefit of repository diversity.

\begin{figure}[t]
\centering
\todo{Scaling plot: x-axis = number of training repos (log scale), y-axis = CR-test EM.
Expected: monotonic improvement with diminishing returns.
Include pretrained baseline as horizontal line.
Generate with: analysis/scaling\_repos.py}
\caption{CR-test EM as a function of training repository count. \method{} benefits from repository diversity, with performance improving log-linearly.}
\label{fig:scaling_repos}
\end{figure}

\todo{Write analysis paragraph:
\begin{itemize}
    \item How many repos are needed to surpass pretrained? FFT?
    \item Is the curve saturating or still improving at 447 repos?
    \item Implications for scaling: collecting more repositories is likely to improve results further.
\end{itemize}}

\subsection{Performance vs.\ Repository Size}

\begin{figure}[t]
\centering
\todo{Scatter plot: x-axis = number of training QnA pairs per repo, y-axis = IR-test EM per repo.
One series per method: pLoRA, Code2LoRA, FFT, Pretrained.
Expected: pLoRA improves with more data per repo; Code2LoRA is more stable.
Generate with: analysis/performance\_vs\_repo\_size.py}
\caption{IR-test EM vs.\ number of training examples per repository. Per-repo LoRA is data-hungry, while \method{} maintains stable performance across repository sizes.}
\label{fig:perf_vs_size}
\end{figure}

\todo{Analyze: pLoRA needs $>$X examples to outperform Code2LoRA. For repos with fewer examples, Code2LoRA's cross-repo transfer is critical.}

\subsection{LoRA Weight Analysis}

We analyze the structure of generated LoRA weights to understand what the hypernetwork learns.

\paragraph{Inter-repository similarity.}
We compute pairwise cosine similarity between flattened LoRA weight vectors generated for different repositories.
\todo{Report:
\begin{itemize}
    \item Average pairwise similarity for Code2LoRA-generated LoRAs.
    \item Average pairwise similarity for independently trained pLoRA adapters.
    \item Expected: Code2LoRA generates more similar adapters (smooth hypernetwork manifold) vs.\ independently trained pLoRA (higher variance).
\end{itemize}}

\begin{figure}[t]
\centering
\todo{t-SNE or UMAP of generated LoRA weights for 50+ repos, colored by repository domain or size.
Generate with: analysis/lora\_similarity.py}
\caption{t-SNE visualization of generated LoRA weight vectors. Repositories with similar codebases produce similar adapters, demonstrating that the hypernetwork captures meaningful repository structure.}
\label{fig:lora_tsne}
\end{figure}

\subsection{Qualitative Examples}

Table~\ref{tab:qualitative} shows representative examples where \method{} succeeds and fails.

\begin{table*}[t]
\centering
\small
\begin{tabular}{p{5cm}p{3cm}p{3cm}p{3cm}}
\toprule
Test Prefix (truncated) & Reference & \method{} & Pretrained \\
\midrule
\todo{Example 1: Code2LoRA correct, pretrained wrong. Pick from CR test---demonstrates cross-repo generalization.} & & & \\
\midrule
\todo{Example 2: Code2LoRA correct, FFT wrong. Shows adaptation advantage over generic fine-tuning.} & & & \\
\midrule
\todo{Example 3: Code2LoRA partially correct (high edit sim). Shows graceful degradation.} & & & \\
\midrule
\todo{Failure: Code2LoRA wrong. Analyze why---rare API, ambiguous target, insufficient repo context.} & & & \\
\bottomrule
\end{tabular}
\caption{Qualitative examples from CR test. \todo{Select from analysis/qualitative\_examples.py output.}}
\label{tab:qualitative}
\end{table*}

\subsection{Effect of Oracle Context Coverage}

Approximately \todo{87.5\%} of QnA pairs have non-empty oracle context (i.e., at least one imported definition was resolved to a repository source file).
The remaining pairs import only from the standard library or third-party packages.

\begin{table}[t]
\centering
\small
\begin{tabular}{lcc}
\toprule
& \multicolumn{2}{c}{CR Test EM (\%)} \\
\cmidrule(lr){2-3}
Method & w/ Oracle & w/o Oracle \\
\midrule
Pretrained & \todoinline{XX.X} & \todoinline{XX.X} \\
Oracle Context & \todoinline{XX.X} & \todoinline{XX.X} \\
\method{} & \todoinline{XX.X} & \todoinline{XX.X} \\
\method{} + Oracle & \todoinline{XX.X} & \todoinline{XX.X} \\
\bottomrule
\end{tabular}
\caption{Performance partitioned by oracle availability on CR test. ``w/ Oracle'' = pairs where oracle context exists; ``w/o Oracle'' = pairs with only stdlib/third-party imports. \todo{Fill from analysis/oracle\_coverage\_analysis.py.}}
\label{tab:oracle_coverage}
\end{table}

\todo{Analyze:
\begin{itemize}
    \item Does oracle context help more for pairs where it's available? (Expected: yes)
    \item Does Code2LoRA still outperform on pairs without oracle? (Expected: yes, showing the embedding captures information beyond import resolution)
\end{itemize}}

\subsection{Efficiency Comparison}

\begin{table}[t]
\centering
\small
\begin{tabular}{lrrl}
\toprule
Method & Extra Tokens & Adapt.\ Time & Storage \\
\midrule
Pretrained & 0 & N/A & 3.0 GB \\
RAG ($k$=3) & $\sim$1,500 & per query & +chunks \\
ICL (3-shot) & $\sim$1,200 & per query & +examples \\
Oracle & $\sim$500 & per query & +cache \\
FFT & 0 & $\sim$\todoinline{X}h & +3.0 GB \\
sLoRA & 0 & $\sim$\todoinline{X}h & +19 MB \\
pLoRA & 0 & $\sim$\todoinline{X}min/repo & +19 MB/repo \\
\method{} & \textbf{0} & $<$10ms/repo & +\todoinline{X} MB \\
\method{}-PAW & \textbf{0} & $<$10ms/repo & +\todoinline{X} MB \\
\bottomrule
\end{tabular}
\caption{Efficiency comparison. \method{} adds zero inference tokens and generates adapters in a single forward pass ($<$10ms on GPU). Storage costs are for the hypernetwork weights only; the frozen base model (3.0 GB) is shared.}
\label{tab:efficiency}
\end{table}

\todo{Measure and fill:
\begin{itemize}
    \item Actual adapter generation time (forward pass through hypernetwork).
    \item Total inference time per sample for each method.
    \item Hypernetwork model sizes in MB.
    \item Training time for FFT, sLoRA.
\end{itemize}}

% ============================================================
\section{Ablation Studies}
\label{sec:ablations}
% ============================================================

We conduct ablation studies on CR val to isolate the contribution of key design choices.

\subsection{LoRA Rank}

\begin{table}[t]
\centering
\small
\begin{tabular}{lcc}
\toprule
LoRA Rank & EM (\%) & Params \\
\midrule
4 & \todoinline{XX.X} & \todoinline{X}M \\
8 & \todoinline{XX.X} & \todoinline{X}M \\
16 (default) & \todoinline{XX.X} & \todoinline{X}M \\
32 & \todoinline{XX.X} & \todoinline{X}M \\
\bottomrule
\end{tabular}
\caption{Effect of LoRA rank on CR val. \todo{Fill from ablation experiments.}}
\label{tab:ablation_rank}
\end{table}

\todo{Analyze: is rank 16 optimal? Does higher rank help or overfit?}

\subsection{Hypernetwork Width}

\begin{table}[t]
\centering
\small
\begin{tabular}{lcc}
\toprule
Hidden Dim & EM (\%) & Params \\
\midrule
256 & \todoinline{XX.X} & \todoinline{X}M \\
512 (default) & \todoinline{XX.X} & \todoinline{X}M \\
1024 & \todoinline{XX.X} & \todoinline{X}M \\
\bottomrule
\end{tabular}
\caption{Effect of hypernetwork hidden dimension on CR val. \todo{Fill from ablation experiments.}}
\label{tab:ablation_width}
\end{table}

\subsection{Number of Bases (PAW)}

\begin{table}[t]
\centering
\small
\begin{tabular}{lcc}
\toprule
Num Bases ($N$) & EM (\%) & Params \\
\midrule
4 & \todoinline{XX.X} & \todoinline{X}M \\
8 & \todoinline{XX.X} & \todoinline{X}M \\
16 (default) & \todoinline{XX.X} & \todoinline{X}M \\
32 & \todoinline{XX.X} & \todoinline{X}M \\
\bottomrule
\end{tabular}
\caption{Effect of number of basis vectors in the PAW-style hypernetwork on CR val. \todo{Fill from ablation experiments.}}
\label{tab:ablation_bases}
\end{table}

\subsection{Embedding Components}

\begin{table}[t]
\centering
\small
\begin{tabular}{lc}
\toprule
Embedding & EM (\%) \\
\midrule
Mean only & \todoinline{XX.X} \\
Max only & \todoinline{XX.X} \\
Mean + Max (default) & \todoinline{XX.X} \\
Mean + Max + Weighted & \todoinline{XX.X} \\
\bottomrule
\end{tabular}
\caption{Effect of repository embedding composition on CR val. \todo{Fill from ablation experiments.}}
\label{tab:ablation_embedding}
\end{table}

\subsection{Training Data Size}

\begin{table}[t]
\centering
\small
\begin{tabular}{lc}
\toprule
Training Repos & CR Test EM (\%) \\
\midrule
50 (12\%) & \todoinline{XX.X} \\
100 (24\%) & \todoinline{XX.X} \\
200 (49\%) & \todoinline{XX.X} \\
409 (100\%) & 63.8 \\
\bottomrule
\end{tabular}
\caption{Effect of training set size on CR test performance. \todo{Fill from scaling experiments.}}
\label{tab:ablation_data}
\end{table}

\subsection{Architecture Comparison}

\begin{table}[t]
\centering
\small
\begin{tabular}{lccc}
\toprule
Variant & CR EM & IR EM & Params \\
\midrule
Direct-Projection & 63.8 & 66.2 & \todoinline{X}M \\
Shared-Basis (PAW) & \todoinline{XX.X} & \todoinline{XX.X} & \todoinline{X}M \\
\bottomrule
\end{tabular}
\caption{Comparison of hypernetwork architectures. \todo{Fill PAW results after training completes.}}
\label{tab:ablation_arch}
\end{table}

\todo{Analyze the trade-off: PAW uses $\sim$8.5$\times$ fewer parameters. If performance is comparable, PAW is strictly preferable. If there is a gap, discuss the parameter-performance trade-off.}

% ============================================================
\section{Scaling Analysis}
\label{sec:scaling}
% ============================================================

We evaluate \method{} across different base model sizes to understand scaling behavior.

\begin{table}[t]
\centering
\small
\begin{tabular}{lccc}
\toprule
Base Model & Pretrained & FFT & \method{} \\
\midrule
Qwen2.5-Coder-0.5B & \todoinline{XX.X} & \todoinline{XX.X} & \todoinline{XX.X} \\
Qwen2.5-Coder-1.5B & 45.7 & 51.4 & 63.8 \\
Qwen2.5-Coder-3B & \todoinline{XX.X} & \todoinline{XX.X} & \todoinline{XX.X} \\
\bottomrule
\end{tabular}
\caption{Scaling behavior across base model sizes (CR test EM \%). \todo{Fill after scale experiments. Key question: does the benefit of Code2LoRA increase or decrease with model size?}}
\label{tab:scaling}
\end{table}

\todo{Analyze:
\begin{itemize}
    \item Does the gap between Code2LoRA and pretrained grow with model size? (interesting if yes)
    \item Does the gap shrink? (suggests larger models already internalize repo knowledge)
    \item How does FFT scale vs.\ Code2LoRA?
\end{itemize}}

% ============================================================
\section{Discussion}
\label{sec:discussion}
% ============================================================

\paragraph{Why does \method{} outperform fine-tuning?}
FFT adapts all model parameters using training data from all repositories, learning a single ``average'' specialization.
\method{} instead generates repo-\emph{specific} adapters, allowing the model to dynamically specialize for each test repository.
This is particularly beneficial for cross-repo evaluation, where the repository was never seen during training---the hypernetwork transfers structural patterns from training repositories to generate effective adapters for new ones.

\paragraph{Why do context-based methods underperform?}
For assertion completion, the target is typically a short expression whose value depends primarily on the immediate test context (variable assignments, function calls in the test body).
Prepending retrieved code or few-shot examples adds tokens that do not directly help predict the assertion value and may shift the model's distribution.
This does not imply that RAG and ICL are generally ineffective for code tasks---they may excel at tasks requiring explicit cross-file knowledge (e.g., API usage, type inference).
We plan to investigate task-specific interactions in future work.
\todo{Verify this explanation with attention analysis or context ablation.}

\paragraph{Limitations of the current evaluation.}
Per-repo LoRA results are based on 10 repositories (99 test pairs).
Expanding to all IR repositories would provide a more representative comparison but requires substantial compute.
Additionally, the single LoRA baseline was initially trained with 4-bit quantization; we retrained with \texttt{bfloat16} for fair comparison, but results should be verified.
\todo{Verify final sLoRA numbers after bf16 retraining.}

% ============================================================
\section{Conclusion}
\label{sec:conclusion}
% ============================================================

We presented \method{}, a hypernetwork approach for generating repository-specific LoRA adapters from code embeddings.
Given a new repository, \method{} produces task-effective adapters in a single forward pass, requiring no retrieval, no additional inference tokens, and no per-repository retraining.
We introduced two hypernetwork architectures---direct-projection and shared-basis---offering a trade-off between representational capacity and parameter efficiency.

Evaluated on \benchmark{}, a benchmark of 560 Python repositories with 73K assertion completion tasks across cross-repo and in-repo settings, \method{} achieves 63.8\% exact match on cross-repo evaluation---substantially outperforming pretrained (45.7\%), fine-tuned (51.4\%), and context-based (RAG: 39.9\%, ICL: 42.2\%) baselines.
Notably, \method{} outperforms even per-repo LoRA training (57.6\%) on in-repo evaluation (66.2\%) despite requiring zero per-repo training data.

Our results demonstrate that distilling repository knowledge into adapter parameters via a hypernetwork is more effective than injecting it through the input context, opening a new direction for efficient, scalable repository-level code understanding.

\todo{Update final numbers and claims after all experiments complete. Mention composable results if available. Polish conclusion.}

% ============================================================
\section*{Limitations}
% ============================================================

\begin{enumerate}[nosep]
    \item \textbf{Task scope.} We evaluate only on Python assertion completion. Generalization to other programming languages, other code tasks (e.g., function completion, bug fixing), and non-code domains remains to be validated.

    \item \textbf{Evaluation metric.} Exact match does not capture functional equivalence---a different but semantically equivalent expression is scored as incorrect. Edit similarity and CodeBLEU partially address this, but execution-based evaluation would be more robust.

    \item \textbf{Repository scale.} Our benchmark uses 560 repositories. While this is sufficient for initial evaluation, larger-scale experiments would strengthen claims about generalization.

    \item \textbf{Hypernetwork overhead.} The hypernetwork adds $\sim$\todoinline{X}M parameters that must be stored alongside the base model. While adapter generation is fast ($<$10ms), the hypernetwork must be loaded into memory.

    \item \textbf{Static embeddings.} Repository embeddings are computed once and do not update as the codebase changes. The composable design addresses this partially through incremental file-level updates, but a fully online adaptation mechanism is left to future work.
    \todo{Verify composable claim is supported by experiments.}
\end{enumerate}

% ============================================================
\section*{Ethics Statement}
% ============================================================

Our work uses publicly available open-source Python repositories with permissive licenses.
We do not train on or generate personally identifiable information.
The assertion completion task evaluates model accuracy on developer-written tests and does not involve generating novel code that could have safety implications.
\todo{Review ethics requirements for EMNLP 2025.}

% ============================================================
\section*{Acknowledgments}
% ============================================================

\todo{Add acknowledgments: compute resources (Alliance Canada), funding, colleagues.}

% ============================================================
\bibliography{references}
\bibliographystyle{acl_natbib}

% ============================================================
% APPENDIX
% ============================================================
\appendix

\section{Additional Results}
\label{app:additional_results}

\subsection{RAG with Different $k$}

\begin{table}[h]
\centering
\small
\begin{tabular}{lccc@{\hskip 12pt}ccc}
\toprule
& \multicolumn{3}{c}{CR Test} & \multicolumn{3}{c}{IR Test} \\
\cmidrule(lr){2-4} \cmidrule(lr){5-7}
$k$ & EM & EditSim & BLEU & EM & EditSim & BLEU \\
\midrule
3 & 39.9 & 0.534 & 0.439 & 42.3 & 0.554 & 0.447 \\
5 & 37.8 & 0.510 & 0.422 & 39.9 & 0.529 & 0.428 \\
10 & 36.4 & 0.491 & 0.406 & 39.4 & 0.518 & 0.420 \\
\bottomrule
\end{tabular}
\caption{RAG results for varying $k$. Performance decreases with more retrieved chunks, suggesting retrieval noise outweighs additional context.}
\label{tab:rag_k}
\end{table}

\subsection{ICL with Different Shot Counts}

\begin{table}[h]
\centering
\small
\begin{tabular}{lccc@{\hskip 12pt}ccc}
\toprule
& \multicolumn{3}{c}{CR Test} & \multicolumn{3}{c}{IR Test} \\
\cmidrule(lr){2-4} \cmidrule(lr){5-7}
Shots & EM & EditSim & BLEU & EM & EditSim & BLEU \\
\midrule
3 & 42.2 & 0.560 & 0.470 & 45.3 & 0.593 & 0.482 \\
5 & 42.0\textsuperscript{*} & 0.544 & 0.462 & 44.9 & 0.587 & 0.479 \\
\bottomrule
\end{tabular}
\caption{ICL results for varying shot counts. \textsuperscript{*}5-shot CR evaluated on n=2,750 (partial run). \todo{Complete 5-shot CR run.}}
\label{tab:icl_shots}
\end{table}

\subsection{Full Per-Repository Breakdown}
\label{app:per_repo}

\todo{Add table with per-repo EM for all IR-test repos across methods. Or top/bottom 10 repos.}

\section{Implementation Details}
\label{app:implementation}

\subsection{Dataset Construction Details}

\paragraph{Test file identification.}
Files are classified as test files if they match any of:
\texttt{test\_*.py}, \texttt{*\_test.py}, or reside in directories named \texttt{tests/}, \texttt{test/}.
Identified test files are moved to a separate \texttt{TEST\_HYPERNET/} directory within each repository, preserving relative paths.

\paragraph{Structured prefix construction.}
Each QnA prefix is constructed as follows:
\begin{enumerate}[nosep]
    \item All import statements from the test file.
    \item The enclosing class definition (if the test is a method).
    \item Helper methods (\texttt{setUp}, \texttt{tearDown}, fixtures).
    \item The test function signature and body up to the assertion cut point.
\end{enumerate}
This structured approach preserves the most informative context while managing token budget.

\paragraph{Quality filters applied.}
\begin{itemize}[nosep]
    \item Targets starting with comma (malformed AST segmentation).
    \item Targets outside function bodies (module-level assertions).
    \item Empty or whitespace-only targets.
    \item Duplicate targets within the same test function.
    \item Targets containing only punctuation or single characters.
\end{itemize}

\subsection{Oracle Context Construction}

\paragraph{Import resolution strategy.}
For each import in the test prefix:
\begin{enumerate}[nosep]
    \item Parse using AST with fallback regex for syntax errors.
    \item Resolve the module to a file path, trying multiple source roots: repository root, \texttt{src/}, \texttt{lib/}, package directories with \texttt{\_\_init\_\_.py}.
    \item For relative imports, resolve relative to the test file's location.
    \item If the imported name is used in the test prefix, extract its definition (function or class) from the source file via AST.
\end{enumerate}

\paragraph{Coverage.}
\todo{Report: X\% of pairs have oracle, average oracle length in tokens, distribution.}

\subsection{Training Details}

\begin{table}[h]
\centering
\small
\begin{tabular}{lcccc}
\toprule
& FFT & sLoRA & \method{} & \method{}-PAW \\
\midrule
LR & 2e-5 & 2e-4 & 1e-4 & 1e-4 \\
Epochs & 3 & 3 & 3 & 3 \\
Max seq len & 2048 & 2048 & 8192 & 8192 \\
Batch size & \todoinline{X} & \todoinline{X} & 1 & 1 \\
Grad accum & \todoinline{X} & \todoinline{X} & \todoinline{X} & \todoinline{X} \\
Precision & bf16 & bf16 & bf16 & bf16 \\
Optimizer & AdamW & AdamW & AdamW & AdamW \\
\bottomrule
\end{tabular}
\caption{Training hyperparameters for all trained methods. \todo{Fill missing values.}}
\label{tab:training_details}
\end{table}

\subsection{Compute Resources}

All experiments were conducted on Alliance Canada HPC clusters using NVIDIA H100 GPUs.
\todo{Report total GPU hours, specific node types, wallclock times for each experiment.}

\section{Broader Analysis}
\label{app:analysis}

\todo{Add additional analysis figures here:
\begin{itemize}
    \item Learning curves (training loss over steps for each method).
    \item Attention pattern visualization showing how LoRA-adapted model differs from pretrained.
    \item Error categorization table (types of failures: wrong value, wrong type, hallucinated API, etc.).
\end{itemize}}

\end{document}
